\originalchapter{真空能、加速膨胀与参数推断}
本章对应8.286第18--21讲，并接回第1--2讲中的真空能问题。背景章给出了含$\Lambda$的演化式和解析年龄，距离章给出了光线测距。本章说明“负压强”的物理含义，以及观测为何需要多种证据共同约束。量子场论的正规化推导见QFT卷；这里不把一个发散的零点能积分当成可直接读出的宇宙常数。

\originalsection{恒定能量密度怎样产生负压强}
取真空能密度$\epsilon_v$不随体积变化，区域能量$U=\epsilon_vV$。若准静态改变体积而无热流，$\dd U=-p\dd V$给$p_v=-\epsilon_v$。假想活塞的外面真空压强取零，拉开活塞需对内部系统做正功；这解释了为何负压强并非“负能量密度”。真实宇宙无这样的活塞边界，但局部共动能量守恒给同一关系。

Lorentz不变真空不能挑选一个静止速度方向，故其应力必须与度规成比例：$T_{\mu\nu}^{v}=-\epsilon_vg_{\mu\nu}$。把它移到Einstein方程的左边可写成
\[
 \Lambda=\frac{8\pi G\epsilon_v}{c^4},\qquad
 \rho_v=\frac{\epsilon_v}{c^2},\quad p_v=-\rho_vc^2.
\]
加速方程中的$\rho+3p/c^2$因此贡献$-2\rho_v$。这仍是正$G$和正能量密度的引力理论，不是把Newton常数改成负数。在真空能主导且展开的阶段，$a\sim\ee^{H_vt}$，$H_v=\sqrt{8\pi G\rho_v/3}$。加入其他组分或曲率时不是从一开始就精确指数。

膨胀时$U$可增大而$\epsilon_v$不变，满足$\dd U=-p_v\dd V$。这不授权宣称整个宇宙有一个定义无歧义且值为零的全局总能量。一般FLRW时空缺少全局时间平移对称性；第1--2讲用负引力能与正物质能抵消作直观图像，本卷将其保留为特定模型和能量定义下的启发，不作为已证明的普遍定理。

\originalsection{由现象到组分：三种证据各约束什么}
远处Ia超新星在相同红移下偏暗，意味着在标准化光度、传播和选择效应受控后光度距离较大。距离是$H(z)$的累计积分，不能由某一个超新星直接测出瞬时$\ddot a$。原第21讲所述20--30\%偏暗是1998年研究的量级描述，具体依赖比较模型和红移，不是任何$z$都保持同一差值。

CMB声学角尺度主要比较最后散射时的声学标准尺与角直径距离，帮助限制空间曲率和组分，但存在几何退化。其他观测约束物质密度后，近空间平坦而物质不足的组合指向额外组分。恒星年龄则限制宇宙年龄不能比其中已形成的恒星更小。三者从距离、早期标准尺和时钟给互补信息；任何一个结果都需说明校准与模型依赖。

用$w(a)=p_{\rm de}/\epsilon_{\rm de}$描述可变化暗能量，独立守恒给
\[
 \rho_{\rm de}(a)=\rho_{{\rm de},0}
 \exp\left[-3\int_1^a(1+w(a'))\dd\ln a'\right].
\]
$w=-1$恢复宇宙学常数；$w$常数给$a^{-3(1+w)}$。正能量组分的$w<-1/3$可以促加速，却不一定就是严格真空。改变引力理论也可能改变从观测到组分的推断。本卷以GR和FLRW为计算前提，不将“暗能量”这个现象名称等同于唯一已确定的微观粒子。

\originalsection{超新星怎样从标准烛光变为实际拟合}
距离模数是$\mu=m-M=5\log_{10}[D_L/(10\ \mathrm{pc})]$。Ia超新星并非天然相同光度，要根据光变宽度与颜色作标准化。例如8.902 PS5采用$\mu=m-M+\alpha s-\beta c$，其$M$定义与写$\alpha(s-1)$的常见记法相差零点。红移$z$、色参数$c$及速度光速$c$用同一个字母时尤其需在语境中辨清。

因为$D_L\propto H_0^{-1}$，若没有外部绝对光度校准，改变$H_0$可由改变$M$抵消。此时超新星主要限制距离曲线\emph{形状}和组分，不能单独测出$H_0$。更完整分析还需校准误差、相关光变参数协方差、选择效应、宿主环境、吸收与模型内禀散布。后面真实SNLS习题已实际读取数据，按题给简化模型拟合，并明确两种散布解释；那是课程计算练习，不冒充发表论文的完整联合似然。

\originalsection{年龄与距离的积分边界不可省略}
对今天$a_0=1$的展开模型，含辐射、物质、曲率、真空能时
\[
 H_0t_0=\int_0^1\frac{a\dd a}
 {\sqrt{\Omega_{r,0}+\Omega_{m,0}a+\Omega_{K,0}a^2+\Omega_{v,0}a^4}}.
\]
这对应第20讲的年龄积分；只有在积分区间内根号为正且存在所选大爆炸展开分支才适用。允许负$\Lambda$、回转或反弹时须先查分支，不能把所有密度参数直接丢进一个数值积分再解释为年龄。

光度距离还含空间面积因子。正曲率时$D_M=R_{K0}\sin(\chi/R_{K0})$，负曲率用$\sinh$，$\chi=c\int_0^z\dd z'/H(z')$。红移光子的能量与到达率各降一个$1+z$，故$D_L=(1+z)D_M$。正曲率反极点附近需处理焦点和可能的多路径，普通近邻资料的模型适用区间与整个封闭空间的光线拓扑不能混成同一个简单距离分支。

\originalsection{为何真空能量级仍是独立的理论难题}
在无引力的实验里，势能加常数不改变相对动力学；有引力时常数进入$T_{\mu\nu}$，影响宇宙膨胀。量子自由场的每个谐振模式有零点能，形式上玻色场贡献$+\hbar\omega/2$，费米场的真空贡献符号不同；Higgs等场在最低势点的常数也参与记账。

以动量截止$\Lambda_{\rm UV}$粗估，一种无质量场的零点能密度随$\Lambda_{\rm UV}^4$增长。截止取Planck能标会与宇宙所需小尺度相差约$10^{120}$量级。这说明参数自然性和不同来源之间的抵消问题；它不等于用不受控截止已精确预测了$\Lambda$。有效场论中还存在允许的宇宙学常数反项，最终可观测的是重整化后的总量，而不是某一单项的发散。

第1--2、21讲讨论多真空与环境选择：若真空能过大而正，结构增长过早停止；若过大而负，宇宙可能过早坍缩。因此观察者所在区域可能有选择偏向。但给出可检验概率仍需真空分布、形成权重、观察者代理量与测度。课程的$10^{500}$景观数量是特定理论讨论中的估计，不是已观测的宇宙数；本卷将该框架和争议保留为进阶方向，不据它宣称已解决真空能问题。
